\documentclass[a4paper]{article}

\usepackage[portuguese]{babel}
\usepackage[utf8]{inputenc}
\usepackage[T1]{fontenc}

\usepackage{float}
\usepackage{graphicx}
\usepackage{tgschola}
\usepackage{microtype}
\linespread{1.1}

\title{MAC0216 --- Técnicas de Programação\\Relatório Técnico do Trabalho Prático 1}
\author{
    Ezequiel Gomes de Carvalho - NUSP \\
    Igor Nascimento Silva - NUSP \\
    Lucas Pires da Silva - NUSP \\
    Miguel da Silva Damasceno - NUSP
}
\date{\textit{\small{\today}}}

\begin{document}

\maketitle

\section{Introdução}
A construção do Cbuild - Uma ferramenta de automatização para projetos em c - foi o foco desse trabalho, feito em script (Bash) com conceitos de versionamento de código, moduralização lógica através de funções e tratametnos de erros.

\section{Estrutura}
A ferraemnta foi feita em um único arquivo, com todos os seus comportamentos divididos em funções. O seu uso é feito atraves da linha de comando do Bash.

\section{Arquitetura}
A arquitetura do script é composta por três camadas principais:
\begin{enumerate}
    \item \textbf{Camada de Configuração (Flags):}
    \item \textbf{Camada de Módulos (Funções):}
    \item \textbf{Camada de Roteamento (Case):}
\end{enumerate}

Essa decisão de projeto nos ajudou no desenvolvimento, pois consegumos delegar cada camada para 1 ou 2 integrantes e trabalharmos separademtne em cada camada, onde um erro não interferia no desenvolvimento do ouro sub-grupo.

\subsection{build}

\subsection{run}

\subsection{clean}

\subsection{rebuild}

\subsection{info}

\end{document}
